%============================================================
%  Conclusion
%============================================================
\chapter*{Conclusion and Future Work}
\addcontentsline{toc}{chapter}{Conclusion and Future Work}
\label{chap:conclusion}

\section*{Project Summary}

This end-of-study project delivered a complete, production-ready Vulnerability Management System addressing the critical gap between vulnerability detection and systematic remediation. The VMS transforms raw Nessus scan data into a managed, auditable remediation workflow with enforced SLA deadlines, AI-assisted guidance, and automated stakeholder notifications.

\section*{Key Accomplishments}

Over the course of this project, the following milestones were achieved:

\begin{enumerate}
  \item \textbf{Complete Architecture}: Designed and implemented a three-tier architecture with a 10-table relational database, RESTful API layer, and responsive web interface
  \item \textbf{Nessus Integration}: Built a robust scan import pipeline handling deduplication via composite key, SLA calculation, and automated finding classification
  \item \textbf{Full Lifecycle Management}: Implemented a six-state finding lifecycle with role-based access control for four user types
  \item \textbf{AI Remediation Assistant}: Deployed Ollama/Mistral 7B locally, achieving 85\% useful response quality with zero data privacy compromise
  \item \textbf{Lab Validation}: Tested against 131 real vulnerabilities across three AlmaLinux VMs, with 100\% import accuracy and all functional requirements passing
  \item \textbf{Performance Targets}: Dashboard loads in 1.23 seconds (P50), well below the 2-second target
\end{enumerate}

\section*{Technical Innovation: Claude API to Ollama Migration}

\begin{importantbox}
\textbf{Key Innovation:} The migration from cloud-based Claude API (claude-haiku-4) to locally deployed Ollama/Mistral 7B represents the most significant technical and ethical decision in this project.

Early prototyping used Claude Haiku for generating remediation suggestions. While the response quality was excellent, transmitting vulnerability details — including CVE identifiers, internal IP addresses, open ports, and service versions — to external cloud servers is fundamentally incompatible with security operations best practices.

The migration to Ollama/Mistral 7B required adaptation of the context-building logic and acceptance of somewhat higher response latency (18 seconds vs 3 seconds for cloud APIs). However, it achieves full data sovereignty: no security-sensitive information ever leaves the organization's network. This design decision aligns the VMS with enterprise security policies and makes it deployable in regulated industries (finance, healthcare, government) where cloud AI usage may be prohibited.
\end{importantbox}

\section*{Business Impact}

The VMS provides measurable improvements over manual vulnerability management:

\begin{table}[H]
\centering
\caption{Business Impact — Before vs After VMS}
\begin{tabular}{lll}
\toprule
\textbf{Metric} & \textbf{Before (Manual)} & \textbf{After (VMS)} \\
\midrule
Mean time to assign finding & 3--5 days & $<$ 1 hour \\
SLA breach detection & Weekly review & Real-time \\
Finding deduplication & Manual/Excel & Automatic \\
Remediation guidance & Google search & Context-aware AI \\
Audit trail & Partial & Complete \\
Overdue notification & Ad-hoc & Automated \\
\bottomrule
\end{tabular}
\end{table}

\section*{Limitations}

\begin{itemize}
  \item \textbf{Ollama Latency}: CPU-based inference averages 18 seconds. Production deployment should use GPU-enabled hardware for 3--5 second response times
  \item \textbf{Nessus Dependency}: The system currently requires Tenable Nessus; future versions could support OpenVAS and other scanners
  \item \textbf{Single-Tenant}: The current architecture is single-tenant; multi-organization deployment would require tenant isolation
  \item \textbf{Manual Scanning}: Nessus scan scheduling is managed externally; automated scan triggering within the VMS would improve workflow continuity
\end{itemize}

\section*{Future Work}

\begin{enumerate}
  \item \textbf{Multi-Scanner Support}: Integrate OpenVAS, Qualys, and Rapid7 alongside Nessus
  \item \textbf{JIRA/ServiceNow Integration}: Auto-create tickets when findings are assigned
  \item \textbf{SIEM Connector}: Push findings and status changes to Elasticsearch/Splunk
  \item \textbf{Threat Intelligence}: Enrich findings with real-time EPSS scores and CISA KEV data
  \item \textbf{GPU Acceleration}: Deploy Ollama with CUDA support for sub-3-second AI responses
  \item \textbf{REST API}: Expose a complete REST API for programmatic access and CI/CD pipeline integration
  \item \textbf{Compliance Reporting}: Add PCI-DSS, ISO 27001, and NIS2 compliance report templates
  \item \textbf{Multi-tenancy}: Support multiple organizations in a single deployment with data isolation
\end{enumerate}

\section*{Lessons Learned}

\begin{itemize}
  \item \textbf{Data privacy in AI}: Cloud APIs and security data are fundamentally incompatible. Local LLMs are the right architecture for security tooling
  \item \textbf{Deduplication is critical}: Without the composite key deduplication strategy, repeat scans would produce thousands of duplicate records
  \item \textbf{SLA enforcement drives behavior}: Once deadlines are visible and notifications are automated, remediation velocity increases significantly
  \item \textbf{Iterative development works}: Building the system chapter by chapter, with functional increments at each sprint, allowed for early validation and course correction
\end{itemize}

\vspace{1em}
\noindent The VMS project demonstrates that a well-designed open-source tool stack — Flask, PostgreSQL, Ollama, and Tenable Nessus — can produce enterprise-grade security tooling without commercial licensing costs. The system is ready for deployment in production environments and forms a solid foundation for the future enhancements described above.
